\chapter[Introdução]{Introdução}

% A Universidade de Brasília (UnB) conta, atualmente, com quatro Campi para a realização de suas atividades: Darcy Ribeiro, UnB Planaltina, UnB Ceilândia e UnB Gama. Estes Campi possuem área construída de cerca de 552.171 m² (CEPLAN, 2017), 12.557 m² (CEPLAN, 2017), 9.827 m² (CEPLAN, 2017) e 6.723 m² (CEPLAN, 2017), respectivamente. São grandes áreas que somam 581.278 m² de construções nas dependências da UnB. Além disso, os campus se localizam em diferentes Regiões Administrativas do Distrito Federal, havendo longas distâncias entre eles. Neste cenário, percebe-se a necessidade de implantação de opções de transporte intra e intercampi dentro da Universidade.
A Universidade de Brasília (UnB) conta, atualmente, com quatro campi para a realização de suas atividades: Darcy Ribeiro, UnB Planaltina, UnB Ceilândia e UnB Gama. Estes campi possuem área construída de cerca de \SI{552 171}{\meter\squared}~\cite{ceplan2017Darcy}, \SI{12 557}{\meter\squared}~\cite{ceplan2017FUP}, \SI{9 827}{\meter\squared}~\cite{ceplan2017FCE} e \SI{6 723}{\meter\squared}~\cite{ceplan2017FGA}, respectivamente. São grandes áreas que somam \SI{581 278}{\meter\squared} de construções nas dependências da UnB. Além disso, os campi se localizam em diferentes Regiões Administrativas do Distrito Federal, havendo longas distâncias entre eles. Neste cenário, percebe-se a necessidade de implantação de opções de transporte intra- e intercampi na Universidade.

% Em 2017, a startup Bird foi a primeira empresa a lançar o serviço de compartilhamento de patinetes elétricos (e-scooter), na Califórnia, Estados Unidos. Desde então, este tipo de serviço foi popularizado e hoje encontra-se disponível em mais de 600 cidades em 53 países (MØLLER et al., 2020).
Em 2017, a \estrangeirismo{startup} Bird foi a primeira empresa a lançar o serviço de compartilhamento de patinetes elétricos (\estrangeirismo{e-scooter}), na Califórnia, Estados Unidos. Desde então, esta categoria de serviço foi popularizado e hoje encontra-se disponível em mais de 600 cidades em 53 países~\cite{moller2020}.

% Um dos principais atrativos da utilização de e-scooters é a possibilidade da diminuição da poluição veicular. De acordo com Boglietti et al. (2021), um carro padrão emite entre 147 e 414 gramas de CO2(dióxido de carbono), enquanto  a utilização de um patinete elétrico provoca entre 64 e 237 gramas, com potencial de redução de 20% a 30% desse valor nos próximos anos. Ademais, Silva (2019) afirma que os patinetes elétricos são tão seguros quanto às bicicletas convencionais.
Um dos principais atrativos da utilização de \estrangeirismo{e-scooters} é a possibilidade da diminuição da poluição veicular. De acordo com \textcite{boglietti2021}, um carro padrão emite entre \qtyrange{147}{414}{\gram} de \ch{CO2} (dióxido de carbono), enquanto  a utilização de um patinete elétrico emite entre \qtyrange{64}{237}{\gram}, com potencial de redução de \qtyrange{20}{30}{\percent} desse valor nos próximos anos. Ademais, \textcite{silva2019} afirma que os patinetes elétricos são tão seguros quanto bicicletas convencionais.

% Em uma pesquisa realizada em 2018 em Indianápolis, Estado Unidos, mais de 450.000 viagens em patinetes elétricos foram analisadas. A pesquisa concluiu que a distância média percorrida nas viagens foi menor que 2 quilômetros e que viagens entre 800 metros e 3,2 quilômetros apresentam maior aproveitamento dos benefícios da utilização dos e-scooters (MATHEW et al., 2019). Assim, os patinetes elétricos mostram-se como uma alternativa adequada para transporte intracampi.
Em uma pesquisa realizada em 2018 em Indianápolis, Estado Unidos, mais de \num{450 000} viagens em patinetes elétricos foram analisadas. A pesquisa concluiu que a distância média percorrida nas viagens foi menor que \SI{2}{\km} e que viagens entre \SI{800}{\meter} a \SI{3,2}{\km} apresentam maior aproveitamento dos benefícios da utilização dos \estrangeirismo{e-scooters}~\cite{matthew2019}. Assim, os patinetes elétricos mostram-se como uma alternativa adequada para transporte intracampi.

% Atualmente, já é oferecido, aos alunos da Universidade de Brasília, o serviço de transporte intercampi. Através dos ônibus com itinerários e horários já definidos, é possível se locomover entre os quatro Campi da instituição de maneira gratuita, mediante apresentação de identificação estudantil (SECOM UNB, 2018).
Atualmente já é oferecido, aos alunos da Universidade de Brasília, o serviço de transporte intercampi. Através dos ônibus com itinerários e horários já definidos, é possível se locomover entre os quatro campi da instituição de maneira gratuita, mediante apresentação de identificação estudantil~\cite{secomunb2018}.

% Para a adoção dos patinetes elétricos como meio de transporte, a legislação vigente no Distrito Federal determina que a sua utilização está restrita a locais de circulação de pedestres (com limite de velocidade de 06 km/h) e ciclovias ou ciclofaixas (com limite de 20 km/h), sendo vetado o seu uso em faixas de rolamento (DETRAN, 2019). Além disso, outras orientações para uma utilização mais segura dos patinetes elétricos foram divulgadas pelo DETRAN-DF (Departamento de Trânsito do Distrito Federal) juntamente ao DER-DF (Departamento de Estradas de Rodagem). Estas incluem recomendações para o correto estacionamento dos patinetes, a utilização, mesmo que não seja obrigatória, de equipamentos de proteção individual e sugestões para a primeira utilização deste meio de transporte, entre outras medidas (DETRAN, 2019).
Para a adoção dos patinetes elétricos como meio de transporte, a legislação vigente no Distrito Federal determina que a sua utilização está restrita a locais de circulação de pedestres (com limite de velocidade de \SI{6}{\km\per\hour}) e ciclovias ou ciclofaixas (com limite de \SI{20}{\km\per\hour}), sendo vetado o seu uso em faixas de rolamento~\cite{detran2019}. Além disso, outras orientações para uma utilização mais segura dos patinetes elétricos foram divulgadas pelo DETRAN\nobreakdash-DF (Departamento de Trânsito do Distrito Federal) juntamente ao DER\nobreakdash-DF (Departamento de Estradas de Rodagem do Distrito Federal). Estas incluem recomendações para o correto estacionamento dos patinetes, a utilização, mesmo que não seja obrigatória, de equipamentos de proteção individual e sugestões para a primeira utilização deste meio de transporte, entre outras medidas~\cite{detran2019}.

% Já o serviço intercampi da UnB, por se tratar de transporte escolar, se submete ao Decreto N° 37332 de 12/05/2016. Neste contexto, são aplicadas, entre outras disposições, normas relacionadas às exigências para a habilitação do condutor, padrões de conservação e caracterização do veículo e as autorizações necessárias para a prestação do serviço (BRASÍLIA, 2016).
Já o serviço intercampi da UnB, por se tratar de transporte escolar, se submete ao Decreto~N°~\num{37332} de 12/05/2016. Neste contexto, são aplicadas, entre outras disposições, normas relacionadas às exigências para a habilitação do condutor, padrões de conservação e caracterização do veículo e as autorizações necessárias para a prestação do serviço~\cite{decreto:37332:2016}.

% Dessa forma, há vários fatores que influenciam na circulação intra- e intercampi, que precisam de uma solução com múltiplas possibilidades, inclusive complementares, que permitam a circulação de alunos e servidores nos 4 campi. Por isso, um ecossistema de serviços que se complementam é a solução que tem grande potencial de abarcar as principais necessidades de mercado no que tange a circulação de pessoas, que é o proposto pelo produto deste projeto. A partir desse produto, as soluções existentes são complementadas, tornando os campi melhor conectados e solucionando a demanda levantada.
Dessa forma, há vários fatores que influenciam na circulação intra- e intercampi, que precisam de uma solução com múltiplas possibilidades, inclusive complementares, que permitam a circulação de alunos e servidores nos 4 campi. Por isso, um ecossistema de serviços que se complementam é a solução com grande potencial de abarcar as principais necessidades de mercado no que tange a circulação de pessoas, que é o proposto pelo produto deste projeto. A partir desse produto, as soluções existentes são complementadas, tornando os campi melhor conectados e solucionando a demanda levantada.
